% Generated draft: requires mathematical and source review.
\documentclass[11pt]{article}
\usepackage[margin=25mm]{geometry}
\usepackage{fontspec}
\setmainfont{Latin Modern Roman}
\usepackage{amsmath,amssymb,mathtools,unicode-math}
\setmathfont{Latin Modern Math}
\usepackage{enumitem,needspace,xurl,hyperref}
\setlength{\parindent}{0pt}
\setlength{\parskip}{5pt}
\setlength{\emergencystretch}{3em}
\newcommand{\N}{\mathbb{N}}
\newcommand{\Z}{\mathbb{Z}}
\newcommand{\Q}{\mathbb{Q}}
\newcommand{\R}{\mathbb{R}}
\newcommand{\C}{\mathbb{C}}
\newcommand{\F}{\mathbb{F}}
\newcommand{\A}{\mathbb{A}}
\newcommand{\PP}{\mathbb{P}}
\newcommand{\E}{\mathbb{E}}
\newcommand{\eps}{\varepsilon}
\newcommand{\dd}{\,\mathrm{d}}
\begin{document}
\section*{UnsolvedMath ID: 10000018}
\textbf{Source problem code:} catalogue entry
\subsection{Definitions and mathematical statement}
In a radius interpretation, let $G$ be a finite connected simple graph on $n$ vertices, with graph distance and diameter $D$. Write $B_G(v,k)$ for the induced radius-$k$ ball with distinguished root $v$. Suppose every pair $u,v$ has a rooted isomorphism
\[
(B_G(u,k),u)\cong(B_G(v,k),v).
\]
The first question asks whether this forces vertex transitivity when $k>n/2$, and separately when $k$ is within a fixed constant of $D$. Vertex transitivity means that for every $u,v$ some automorphism of the whole graph sends $u$ to $v$.

The source says balls of size $k$, not explicitly radius $k$; cardinality and radius are different. It also does not state connectedness or the exact second-threshold quantifier. For a radius parameter, $k\geq D$ makes every rooted ball the whole graph, so the substantive near-diameter range is below $D$. Whether the claim concerns each fixed positive constant $C$, a particular $C$, or existence of one in a bound such as $D-C\leq k<D$ is missing. No such choice is silently made.

The second, separate question samples a vertex-transitive graph $G_n$ uniformly on an odd number $n$ of vertices and asks whether its common degree $d(G_n)$ is concentrated near $(n-1)/2$. A possible relative-concentration formulation is
\[
\forall\varepsilon>0,\qquad
\Pr\left(\left|d(G_n)-\frac{n-1}{2}\right|>\varepsilon n\right)\longrightarrow0
\quad(n\to\infty\text{ through odd integers}).
\]
This is an indicated interpretation only. Uniform sampling of labelled graphs on $\{1,\ldots,n\}$ differs from uniform sampling of isomorphism classes, and conditioning on connectedness changes the model. The source fixes none of these or the concentration scale. An $o(n)$ statement, bounded fluctuations and a specified limiting fluctuation law are distinct targets.
\subsection{Short English statement}
Do identical large rooted neighborhoods force global vertex transitivity? Separately, does a uniformly sampled odd-order vertex-transitive graph have degree close to half its order? Clarify ball size, threshold quantifiers and the random-graph sampling model.
\subsection{Sources}
\begin{itemize}
\item Local UnsolvedMath registry, permanent ID 10000018, exact source code catalogue entry, ``Large identical neighborhoods and vertex transitivity''; attributed to Itai Benjamini, no proposed year supplied. Registry status open, classification OPEN-TRIAGE.
\item Itai Benjamini, \emph{Coarse Geometry and Randomness}, Saint-Flour notes, Open problem 2.7, PDF page 23, \url{https://www.wisdom.weizmann.ac.il/~itai/stflouraug24.pdf#page=23}, preserved at \url{https://arquivo.pt/noFrame/replay/20201231041548id_/http://www.wisdom.weizmann.ac.il/~itai/stflouraug24.pdf}. Supplied pypdf-text extraction checked against a rendered PDF, accessed July 29, 2026. No PDF or rendering was accessed here.
\item The supplied summary reports no verified resolution of either part. Provenance says the item survived cross-document deduplication/source-side review, with no source-side solution and no complete primary result found in its bounded modern audit.
\item Current primary-literature status and release-reuse terms remain NEEDS\_REVIEW.
\end{itemize}
\end{document}
